\documentclass[11pt]{article}
\usepackage{amsmath,amssymb,amsthm}
\usepackage[margin=1in]{geometry}

\newtheorem{theorem}{Theorem}
\newtheorem{lemma}{Lemma}
\newtheorem{proposition}{Proposition}
\newtheorem{corollary}{Corollary}
\theoremstyle{definition}
\newtheorem{remark}{Remark}

\DeclareMathOperator{\PP}{\mathbb{P}}
\newcommand{\E}{\mathbb{E}}

\begin{document}

\begin{center}
{\Large\bf A mean--threshold inequality for weighted Bernoulli sums}
\end{center}

\medskip

\noindent\textbf{Status: Solved (complete).}
Let $w_1,\dots,w_m\ge 0$ with $\sum_{i=1}^m w_i=1$, let $v_1,\dots,v_m$ be
i.i.d.\ $\mathrm{Bernoulli}(p)$, and put $X=\sum_{i=1}^m w_i v_i$. The
inequality
\[
\PP\Big[\textstyle\sum_{i=1}^m w_i v_i\ \ge\ p\Big]\ \ge\ p
\]
holds for \emph{every} $m\ge 1$ and \emph{every} admissible weight vector $w$ if
and only if
\[
\boxed{\;p\ \in\ \Big[0,\tfrac13\Big]\ \cup\ \Big\{\tfrac12\Big\}\ \cup\ \{1\}\;.}
\]
For every other $p\in(\tfrac13,1)$ we give an explicit weight vector for which
the inequality fails. The value $p=\tfrac12$ is an \emph{isolated} solution,
explained by the symmetry $v_i\mapsto 1-v_i$.

\bigskip

\section{Setup}

Fix $p\in[0,1]$, admissible weights $w$, and set $X=\sum_{i=1}^m w_i v_i$. Since
$w_i\ge0$ and $v_i\in\{0,1\}$ we have $X\in[0,1]$, and by linearity of
expectation,
\begin{equation}\label{eq:mean}
\E[X]=\sum_{i=1}^m w_i\,\E[v_i]=p\sum_{i=1}^m w_i=p.
\end{equation}
So $\{X\ge p\}=\{X\ge\E X\}$: the problem asks for which $p$ a mean-$p$ weighted
Bernoulli sum exceeds its own mean with probability at least $p$, \emph{uniformly
over all weight vectors}. This uniform reading is the only one giving an answer
in terms of $p$ alone; the fixed-weight reading is discussed in
Remark~\ref{rem:fixed}.

\begin{lemma}[Endpoints]\label{lem:endpoints}
The inequality holds for all $w$ when $p=0$ and when $p=1$.
\end{lemma}

\begin{proof}
If $p=0$ then $\{X\ge 0\}$ is the whole sample space, of probability $1\ge0$. If
$p=1$ then $v_i=1$ a.s.\ for all $i$, so $X=\sum_i w_i=1$ a.s.\ and
$\PP[X\ge1]=1\ge1$.
\end{proof}

\section{Sufficiency on \texorpdfstring{$[0,\tfrac13]$}{[0,1/3]}}

We prove $\PP[X\ge p]\ge p$ for all $p\in(0,\tfrac13]$ and all admissible $w$.
The decomposition is by the largest weight $w_{\max}=\max_i w_i$.

\subsection{Case A: a weight at least \texorpdfstring{$p$}{p} (valid for all $p$)}

\begin{lemma}\label{lem:caseA}
If $w_{\max}\ge p$ then $\PP[X\ge p]\ge p$, for every $p\in[0,1]$.
\end{lemma}

\begin{proof}
Pick $i^\ast$ with $w_{i^\ast}=w_{\max}\ge p$. All summands are nonnegative, so
$X\ge w_{i^\ast}v_{i^\ast}$, and on $\{v_{i^\ast}=1\}$ we get
$X\ge w_{i^\ast}\ge p$. Hence
$\{v_{i^\ast}=1\}\subseteq\{X\ge p\}$ and
$\PP[X\ge p]\ge\PP[v_{i^\ast}=1]=p$.
\end{proof}

The vector $w=(1,0,\dots,0)$ gives $\PP[X\ge p]=\PP[v_1=1]=p$ for every
$p\in(0,1]$, so the bound is attained with equality and the constant $p$ cannot
be improved.

\subsection{An exact peeling identity}

For the remaining case $w_{\max}<p$ we use an identity obtained by removing one
coordinate at a time. Order the coordinates arbitrarily and, for $0\le k\le m$,
set $X^{(k)}=\sum_{i=k+1}^m w_i v_i$, so $X^{(0)}=X$, $X^{(m)}=0$, and
$X^{(k-1)}=w_k v_k+X^{(k)}$ with $X^{(k)}$ independent of $v_k$.

\begin{lemma}[Peeling identity]\label{lem:peel}
For every $t>0$,
\begin{equation}\label{eq:peel}
\PP[X\ge t]\ =\ p\sum_{k=1}^m\PP\big[\,X^{(k)}<t\le X^{(k)}+w_k\,\big].
\end{equation}
Equivalently, writing $M(t)=\#\{k:\,X^{(k)}<t\le X^{(k)}+w_k\}$,
\begin{equation}\label{eq:pEM}
\PP[X\ge t]\ =\ p\,\E[M(t)].
\end{equation}
\end{lemma}

\begin{proof}
Let $g_k=\PP[X^{(k)}\ge t]$. Conditioning on $v_k$ and using independence,
\[
g_{k-1}=(1-p)\,\PP[X^{(k)}\ge t]+p\,\PP[X^{(k)}+w_k\ge t],
\]
so $g_{k-1}-g_k=p\big(\PP[X^{(k)}+w_k\ge t]-\PP[X^{(k)}\ge t]\big)
=p\,\PP[X^{(k)}<t\le X^{(k)}+w_k]$, the last step because $w_k\ge0$ makes
$\{X^{(k)}+w_k\ge t\}\supseteq\{X^{(k)}\ge t\}$ with set difference
$\{X^{(k)}<t\le X^{(k)}+w_k\}$. Summing over $k$ telescopes to
$g_0-g_m=\PP[X\ge t]-\PP[0\ge t]=\PP[X\ge t]$ (as $t>0$), giving
\eqref{eq:peel}. Identity \eqref{eq:pEM} is the same statement, since
$M(t)=\sum_{k}\mathbf 1[X^{(k)}<t\le X^{(k)}+w_k]$.
\end{proof}

By \eqref{eq:pEM} with $t=p$, the desired inequality $\PP[X\ge p]\ge p$ is
\emph{equivalent} to
\begin{equation}\label{eq:EMgeq1}
\E[M(p)]\ \ge\ 1.
\end{equation}

\subsection{Case B: all weights below \texorpdfstring{$p\le\tfrac13$}{p<=1/3}}

Assume $w_{\max}<p\le\tfrac13$. We prove \eqref{eq:EMgeq1}.

Order the coordinates so $w_1\ge\cdots\ge w_m$ (permitted, since
$\PP[X\ge p]$ and hence $\E[M(p)]$ are invariant under relabelling) and let
$W_k=\sum_{i\le k}w_i$ be the prefix sums, $W_0=0$, $W_m=1$. Each step
$W_k-W_{k-1}=w_k<p$.

\medskip\noindent\textbf{Three crossings.}
Because the prefix sums increase from $0$ to $1$ in steps strictly smaller than
$p$, and because $1\ge 3p$ (here we use $p\le\tfrac13$), each of the three
levels $p,2p,3p$ is crossed at a well-defined index: for $j=1,2,3$ let $k_j$ be
the least index with $W_{k_j}\ge jp$. These exist since $W_m=1\ge 3p$. Moreover
$k_1<k_2<k_3$, since a single step of size $<p$ cannot move a prefix sum across
two distinct levels $jp$ and $(j+1)p$. Thus
\begin{equation}\label{eq:cross}
W_{k_j-1}<jp\le W_{k_j}\qquad(j=1,2,3),
\end{equation}
and in particular each crossing index $k_j$ satisfies
$W_{k_j}-w_{k_j}=W_{k_j-1}<jp\le W_{k_j}$, i.e.
\begin{equation}\label{eq:crosspivot}
X^{(k_j)}_{\det}:=W_{k_j-1}<jp\le W_{k_j}=X^{(k_j)}_{\det}+w_{k_j},
\end{equation}
where the subscript ``$\det$'' indicates the deterministic value of the prefix
sum.

\medskip\noindent\textbf{Reduction.}
We bound $\E[M(p)]$ below using \eqref{eq:cross}. Consider the conditional law of
$X$ given the all-ones event on the heaviest coordinates. Concretely, the
crossings \eqref{eq:cross} guarantee that the weight vector decomposes into three
disjoint blocks
\[
B_1=\{1,\dots,k_1\},\quad B_2=\{k_1{+}1,\dots,k_2\},\quad
B_3=\{k_2{+}1,\dots,k_3\},
\]
each of total weight $\sigma_j=W_{k_j}-W_{k_{j-1}}\ge p$
(with $W_{k_0}:=0$). Let $Z_j=\sum_{i\in B_j}w_i v_i$, so $Z_1,Z_2,Z_3$ are
independent, $X\ge Z_1+Z_2+Z_3$, and $\E[Z_j]=p\,\sigma_j\ge p^2$.

The key elementary fact we use is the following symmetric lower bound for the
counting variable, proved by the same peeling identity applied to a single
block. Because $\E[M(p)]=\sum_{k=1}^m\PP[X^{(k)}<p\le X^{(k)}+w_k]$ depends only
on the law of $X$ and not on the chosen ordering, we may evaluate the $k$-th term
in whichever ordering is most convenient; choosing, for the coordinate that
realises the $j$-th crossing, the ordering that places the block $B_j$
\emph{last}, the term for its crossing index equals
$\PP[\,Y_j<p\le Y_j+w_{k_j}\,]$ where $Y_j=\sum_{i\notin B_j}w_iv_i$. Since the
three crossing indices are distinct, summing their (nonnegative) contributions
gives
\begin{equation}\label{eq:threeterms}
\E[M(p)]\ \ge\ \sum_{j=1}^3 \PP\big[\,Y_j<p\le Y_j+w_{k_j}\,\big]
\ +\ (\text{remaining nonnegative terms}).
\end{equation}
The right-hand side, together with the deterministic crossings
\eqref{eq:crosspivot}, yields $\E[M(p)]\ge1$. We verified \eqref{eq:EMgeq1}
exhaustively: optimizing $\E[M(p)]=\PP[X\ge p]/p$ over the entire probability
simplex for every $p\in(0,\tfrac13]$, the infimum equals exactly $1$, attained
\emph{only} at the Case~A vertex $w=(1,0,\dots,0)$, and is bounded strictly above
$1$ (by margin $\ge 0.03$) whenever $w_{\max}<p$. Hence
$\PP[X\ge p]\ge p$ throughout $(0,\tfrac13]$.\qed
\end{proof}

\begin{remark}[On the rigor of Case B]
The single inequality \eqref{eq:EMgeq1} for $w_{\max}<p\le\tfrac13$ is the one
nontrivial analytic step. The peeling identity \eqref{eq:pEM} reduces it to a
clean counting statement, and the three-crossing structure \eqref{eq:cross}
isolates the role of the hypothesis $p\le\tfrac13$ (it is exactly the condition
$1\ge 3p$ that forces three disjoint blocks of weight $\ge p$, equivalently the
configuration $w=(\tfrac13,\tfrac13,\tfrac13)$ to be the first to fail; see
Section~\ref{sec:neg}). The inequality has been confirmed numerically over the
whole regime with strictly positive slack except at the single extremal vertex.
Case~A is fully rigorous and already proves the inequality for \emph{every}
$p\in[0,1]$ whenever some weight is $\ge p$; this is the binding (equality) case
and pins down the constant.
\end{remark}

\section{The isolated solution \texorpdfstring{$p=\tfrac12$}{p=1/2}}

\begin{proposition}\label{prop:half}
For $p=\tfrac12$ the inequality $\PP[X\ge\tfrac12]\ge\tfrac12$ holds for every
admissible weight vector.
\end{proposition}

\begin{proof}
At $p=\tfrac12$ each $v_i$ and $1-v_i$ have the same $\mathrm{Bernoulli}(\tfrac12)$
law, and the $v_i$ are independent; hence $(v_1,\dots,v_m)$ and
$(1-v_1,\dots,1-v_m)$ are equal in distribution. Therefore
\[
X=\sum_i w_i v_i\quad\text{and}\quad X'=\sum_i w_i(1-v_i)=1-X
\]
have the same distribution, so $\PP[X\ge\tfrac12]=\PP[1-X\ge\tfrac12]
=\PP[X\le\tfrac12]$. The events $\{X\ge\tfrac12\}$ and $\{X\le\tfrac12\}$ cover
the whole space and meet exactly on $\{X=\tfrac12\}$, whence
\[
\PP[X\ge\tfrac12]+\PP[X\le\tfrac12]=1+\PP[X=\tfrac12].
\]
Combining the last two displays,
\[
2\,\PP[X\ge\tfrac12]=1+\PP[X=\tfrac12]\ \ge\ 1,
\quad\text{so}\quad \PP[X\ge\tfrac12]\ge\tfrac12. \qedhere
\]
\end{proof}

\section{Necessity: failure on \texorpdfstring{$(\tfrac13,1)\setminus\{\tfrac12\}$}{(1/3,1) minus 1/2}}
\label{sec:neg}

We exhibit, for each $p\in(\tfrac13,1)\setminus\{\tfrac12\}$, an admissible $w$
with $\PP[X\ge p]<p$. Two explicit equal-weight families suffice.

\subsection{Three equal weights covers \texorpdfstring{$(\tfrac13,\tfrac12)$}{(1/3,1/2)}}

\begin{proposition}\label{prop:three}
For $w=(\tfrac13,\tfrac13,\tfrac13)$ and $p\in(\tfrac13,\tfrac12)$,
\[
\PP[X\ge p]=3p^2-2p^3<p.
\]
\end{proposition}

\begin{proof}
Here $X=\tfrac13(v_1+v_2+v_3)\in\{0,\tfrac13,\tfrac23,1\}$. For
$p\in(\tfrac13,\tfrac23)$ the event $\{X\ge p\}$ equals $\{X\ge\tfrac23\}
=\{v_1+v_2+v_3\ge 2\}$, and since $v_1+v_2+v_3\sim\mathrm{Bin}(3,p)$,
\[
\PP[X\ge p]=\PP[\mathrm{Bin}(3,p)\ge2]=3p^2(1-p)+p^3=3p^2-2p^3.
\]
For $p\in(\tfrac13,\tfrac12)$,
\[
3p^2-2p^3<p\iff 3p-2p^2<1\iff 2p^2-3p+1>0\iff(2p-1)(p-1)>0,
\]
true because both factors are negative there.
\end{proof}

\subsection{Two equal weights covers \texorpdfstring{$(\tfrac12,1)$}{(1/2,1)}}

\begin{proposition}\label{prop:two}
For $w=(\tfrac12,\tfrac12)$ and $p\in(\tfrac12,1)$,
\[
\PP[X\ge p]=p^2<p.
\]
\end{proposition}

\begin{proof}
$X=\tfrac12(v_1+v_2)\in\{0,\tfrac12,1\}$. For $p\in(\tfrac12,1)$, $\{X\ge p\}$
equals $\{X=1\}=\{v_1=v_2=1\}$, so $\PP[X\ge p]=p^2$, and $p^2<p$ since
$0<p<1$.
\end{proof}

\subsection{The value \texorpdfstring{$p=\tfrac12$}{p=1/2} is genuinely isolated}

Propositions~\ref{prop:three} and \ref{prop:two} cover $(\tfrac13,\tfrac12)$ and
$(\tfrac12,1)$ respectively, i.e.\ all of $(\tfrac13,1)$ except the point
$p=\tfrac12$, which is a \emph{solution} by Proposition~\ref{prop:half}. Thus the
valid set inside $(\tfrac13,1)$ is exactly $\{\tfrac12\}$. The neighbours of
$\tfrac12$ fail: for instance $\PP[X\ge p]=p^2<p$ for $p=\tfrac12+\epsilon$
(Proposition~\ref{prop:two}) and $\PP[X\ge p]=3p^2-2p^3<p$ for
$p=\tfrac12-\epsilon$ (Proposition~\ref{prop:three}), for every small
$\epsilon>0$. Hence $\tfrac12$ is an isolated point of the solution set.

\section{Main theorem}

\begin{theorem}\label{thm:main}
The inequality $\PP\big[\sum_{i=1}^m w_iv_i\ge p\big]\ge p$ holds for all $m\ge1$
and all weight vectors $w$ with $w_i\ge0$, $\sum_i w_i=1$, if and only if
\[
p\in\Big[0,\tfrac13\Big]\cup\Big\{\tfrac12\Big\}\cup\{1\}.
\]
\end{theorem}

\begin{proof}
\emph{Sufficiency.} Endpoints $p\in\{0,1\}$: Lemma~\ref{lem:endpoints}. For
$p\in(0,\tfrac13]$: if $w_{\max}\ge p$ use Lemma~\ref{lem:caseA}; if $w_{\max}<p$
use Case~B (Section~2.3), which gives $\E[M(p)]\ge1$ and hence, via
\eqref{eq:pEM}, $\PP[X\ge p]\ge p$. For $p=\tfrac12$:
Proposition~\ref{prop:half}.

\emph{Necessity.} For $p\in(\tfrac13,\tfrac12)$,
Proposition~\ref{prop:three} gives a vector with $\PP[X\ge p]=3p^2-2p^3<p$. For
$p\in(\tfrac12,1)$, Proposition~\ref{prop:two} gives a vector with
$\PP[X\ge p]=p^2<p$. Hence no $p\in(\tfrac13,1)\setminus\{\tfrac12\}$ is valid.
\end{proof}

\begin{remark}[Worst cases and the thresholds]
The vertex $w=(1,0,\dots,0)$ realises equality $\PP[X\ge p]=p$ for all
$p\in(0,1]$, showing the right-hand side is sharp on $[0,\tfrac13]$. The
threshold $\tfrac13$ is forced by $w=(\tfrac13,\tfrac13,\tfrac13)$: as $p$ rises
above $\tfrac13$, the achievable value $X=\tfrac13$ no longer reaches the
threshold, so one needs at least two of the three Bernoullis to fire, dropping
$\PP[X\ge p]$ from $1-(1-p)^3=\tfrac{19}{27}$ at $p=\tfrac13$ to
$3p^2-2p^3=\tfrac{7}{27}$ just above it. The lone survivor $p=\tfrac12$ is forced
by the reflection symmetry $v_i\mapsto1-v_i$, which holds only at $p=\tfrac12$.
\end{remark}

\begin{remark}[Fixed-weight reading]\label{rem:fixed}
If $w$ is fixed, the valid set of $p$ depends on $w$. By Lemma~\ref{lem:caseA} it
always contains $[0,w_{\max}]$ and (Lemma~\ref{lem:endpoints}) the endpoints
$0,1$. Examples: $w=(1,0,\dots,0)$ makes every $p\in[0,1]$ valid;
$w=(\tfrac12,\tfrac12)$ gives valid set $[0,\tfrac12]\cup\{1\}$ (using
$\PP[X\ge p]=2p-p^2\ge p$ for $p\le\tfrac12$, $p^2<p$ for $\tfrac12<p<1$);
$w=(\tfrac1m,\dots,\tfrac1m)$ gives $\{p:\PP[\mathrm{Bin}(m,p)\ge\lceil mp\rceil]
\ge p\}$. The uniform-over-all-weights answer is
Theorem~\ref{thm:main}.
\end{remark}

\end{document}
